\documentclass[11pt]{article}
\usepackage[a4paper,margin=1in]{geometry}
\usepackage[T1]{fontenc}
\usepackage[utf8]{inputenc}
\usepackage{amsmath,amssymb,amsthm,amsfonts,mathtools}
\usepackage{booktabs}
\usepackage{array}
\usepackage{enumitem}
\usepackage{hyperref}

\newtheorem{theorem}{Theorem}
\newtheorem{lemma}{Lemma}
\newtheorem{proposition}{Proposition}
\newtheorem{corollary}{Corollary}
\theoremstyle{definition}
\newtheorem{definition}{Definition}
\newtheorem{remark}{Remark}

\newcommand{\U}{\mathcal{U}_n}
\newcommand{\Ic}{I_c}
\newcommand{\J}{\mathcal{J}}
\newcommand{\Band}[2]{\mathcal{B}_{#1}^{#2}}
\newcommand{\Parity}[2]{\mathcal{P}_{#1}^{#2}}
\newcommand{\Weight}[2]{\mathcal{H}_{#1}^{#2}}
\newcommand{\bits}{\{0,1\}}

\title{Causal Index-Set Calculus for Boolean Networks:\\ Exact Gate Formulae, Ordering Invariance, and Deterministic Reconstruction}
\author{}
\date{\today}

\begin{document}
\maketitle

\begin{abstract}
We develop a causality-first formalism for Boolean networks in which local rules are represented by exact index sets over ordered exhaustive repertoires. The basic object is not a sampled trajectory nor a statistical summary, but the set of input indices at which a gate is active under a declared ordering convention. This yields a deterministic calculus in which monotone, parity, implication, threshold, and canalising gates admit closed-form one-sets; MSB-first and LSB-first enumerations are related by an explicit bit-reversal involution; and mixed-gate network outputs are reconstructed exactly by column-wise composition of local index sets. The method provides a mathematically transparent alternative to brute-force repertoire enumeration and to purely probabilistic learning. Existing validation artefacts support exact equality to exhaustive baselines at medium sizes, exact sampled equality at larger sizes, and clear constant-factor gains for analytic prediction over gate-dispatch or truth-table lookup. The paper is written as a mathematical-computational contribution: biology motivates later applications, but the central results concern causal representation, exact reconstruction, and computational complexity.
\end{abstract}

\section{Introduction}
The central difficulty in Boolean-network analysis is not lack of computational power; it is lack of causal compression. One can always enumerate a repertoire, print a truth table, or run a simulation. What remains elusive is a representation that is at once exact, compositional, and explanatory. For a high-level scientific treatment, this distinction matters: a method that merely reproduces outputs is weaker than a method that identifies the rule structure responsible for them.

This paper advances the following thesis. Ordered exhaustive repertoires contain a non-trivial algebraic structure, and this structure can be exploited to replace brute-force enumeration by exact causal formulae. Once the ordering convention is fixed, the question ``what does this network do?'' can be recast as ``at which repertoire indices is each local rule active, and how do those index sets compose?'' The result is an index-set calculus for Boolean networks.

The perspective is causality-first in three senses. First, it is rule-first rather than correlation-first: the object of study is the generative logic of the system, not only its observed frequencies. Second, it is deterministic rather than statistical: the method aims at exact reconstruction under known dynamics, not at approximate inference under unknown mechanisms. Third, it is compositional: global behaviour is built from local causal rules rather than from opaque end-to-end approximations.

This places the work in dialogue with two important traditions without collapsing into either of them. The first is the tradition of simple rules generating rich behaviour, exemplified by Wolfram's programme on discrete systems. The second is the algorithmic-causality tradition associated with perturbation and generative explanation, developed in part by Zenil and collaborators. Our contribution is more specific: we derive a concrete analytic calculus for ordered Boolean repertoires, prove ordering-control results, and connect the formalism to exact network reconstruction and explicit complexity comparisons.

The paper makes six principal claims.
\begin{enumerate}[label=\textbf{C\arabic*.}]
\item Ordered exhaustive repertoires can be treated as algebraic objects rather than as passive lookup tables.
\item Canonical gate families admit exact, network-aware one-set formulae over those repertoires.
\item Ordering changes are controlled by an explicit involution, so exactness is invariant under representation.
\item Mixed-gate synchronous networks can be reconstructed exactly by composing local index sets.
\item The method is mathematically richer than naive vectorised shortcuts and empirically exact where those shortcuts fail.
\item Complexity improvements arise not by magic removal of the \(2^n\) state space, but by replacing brute-force gate evaluation with explicit causal formulae, by enabling exact query answering without full matrix materialisation, and by reducing the validation burden through principled sampling.
\end{enumerate}

\section{Formal Setup}
Let \(cm \in \bits^{n \times n}\) be the connectivity matrix of a Boolean network, where \(cm_{ij}=1\) means that node \(j\) feeds node \(i\). Let
\[
dynamic=(d_1,\dots,d_n)
\]
be the list of local gate labels, and let
\[
N_i=\{j\in\{1,\dots,n\} : cm_{ij}=1\}
\]
be the connected-input set of node \(i\). Each node has a local map
\[
g_i:\bits^{|N_i|}\to\bits,
\]
and the synchronous one-step network update is
\[
F(x)=\bigl(g_1(x_{N_1}),\dots,g_n(x_{N_n})\bigr), \qquad x\in\bits^n.
\]

\begin{definition}[Ordered Exhaustive Repertoire]
The ordered exhaustive repertoire of size \(n\) is the set \(\bits^n\) together with a declared enumeration of all \(2^n\) binary vectors. We write \(\U=\{1,\dots,2^n\}\) for the corresponding index universe.
\end{definition}

Two orderings are operationally relevant throughout the project.
\begin{itemize}
\item \textbf{LSB-first ordering:} the index \(j\) is associated with \(v(j)=\mathrm{Reverse}(\mathrm{IntegerDigits}(j-1,2,n))\), with weights \(2^{i-1}\).
\item \textbf{MSB-first ordering:} the index \(j\) is associated with \(v(j)=\mathrm{IntegerDigits}(j-1,2,n)\), with weights \(2^{n-i}\).
\end{itemize}

\begin{definition}[Bit-Reversal Involution]
For \(j\in\U\), define
\[
\varphi(j)=1+\mathrm{FromDigits}\Bigl(\mathrm{Reverse}\bigl(\mathrm{IntegerDigits}(j-1,2,n)\bigr),2\Bigr).
\]
Then \(\varphi:\U\to\U\) is an involution.
\end{definition}

The involution \(\varphi\) is not a cosmetic convenience. It is the formal device that makes ordering control explicit. Without it, one risks confusing genuine mathematical differences with mere encoding differences.

\section{Causal Index Sets}
The method begins with a structural observation: each coordinate of an ordered exhaustive repertoire forms periodic bands, and Boolean gates can be expressed as exact set operations on those bands.

\begin{definition}[Bit Bands]
For a node index \(i\in\{1,\dots,n\}\) and a bit value \(a\in\bits\), define
\[
\Band{i}{a}=\{j\in\U : v_i(j)=a\}.
\]
\end{definition}

\begin{definition}[Parity and Weight Classes]
For a connected-input set \(\Ic\subseteq\{1,\dots,n\}\), define
\[
\Parity{\Ic}{1}=\Bigl\{j\in\U : \sum_{i\in\Ic} v_i(j)\equiv 1 \pmod 2 \Bigr\},
\]
\[
\Parity{\Ic}{0}=\Bigl\{j\in\U : \sum_{i\in\Ic} v_i(j)\equiv 0 \pmod 2 \Bigr\},
\]
and for \(r\in\{0,\dots,|\Ic|\}\),
\[
\Weight{\Ic}{r}=\Bigl\{j\in\U : \sum_{i\in\Ic} v_i(j)=r \Bigr\}.
\]
\end{definition}

These three families already encode most of the method. Bands capture fixed-bit conditions, parity classes capture XOR-like structure, and Hamming-weight strata capture threshold logic. What remains is to show that standard gate families are simple set expressions built from them.

\section{Exact Gate Formulae}
The following proposition collects the exact one-set formulae for the gate families currently supported by the repository-level analysis series.

\begin{proposition}[Gate-Family One-Sets]
Let node \(k\) have connected-input set \(\Ic\). Under a fixed repertoire ordering, the set of indices at which node \(k\) outputs one is given as follows.
\begin{center}
\small
\begin{tabular}{p{0.20\linewidth}p{0.67\linewidth}}
\toprule
\textbf{Gate family} & \textbf{Exact one-set formula} \\
\midrule
AND & \(\J_k^{\mathrm{AND}}=\bigcap_{i\in\Ic}\Band{i}{1}\) \\
OR & \(\J_k^{\mathrm{OR}}=\bigcup_{i\in\Ic}\Band{i}{1}\) \\
NAND & \(\J_k^{\mathrm{NAND}}=\U\setminus \J_k^{\mathrm{AND}}\) \\
NOR & \(\J_k^{\mathrm{NOR}}=\bigcap_{i\in\Ic}\Band{i}{0}=\U\setminus \J_k^{\mathrm{OR}}\) \\
XOR & \(\J_k^{\mathrm{XOR}}=\Parity{\Ic}{1}\) \\
XNOR & \(\J_k^{\mathrm{XNOR}}=\Parity{\Ic}{0}\) \\
NOT on input \(r\) & \(\J_k^{\mathrm{NOT}}=\Band{r}{0}\) \\
IMPLIES on pair \((a,b)\) & \(\J_k^{\mathrm{IMP}}=\Band{a}{0}\cup\bigl(\Band{a}{1}\cap\Band{b}{1}\bigr)\), equivalently \(\U\setminus\bigl(\Band{a}{1}\cap\Band{b}{0}\bigr)\) \\
NIMPLIES on pair \((a,b)\) & \(\J_k^{\mathrm{NIMP}}=\Band{a}{1}\cap\Band{b}{0}\) \\
KOFN threshold & \(\J_k^{\mathrm{KOFN}}=\bigcup_{r=k}^{|\Ic|}\Weight{\Ic}{r}\) \\
Strict-EXACT-\(k\) & \(\J_k^{=k}=\Weight{\Ic}{k}\) \\
Canalising with trigger set \(C\) and fallback one-set \(\J_h\) & \(\J_k^{\mathrm{CAN}}=C\cup\bigl((\U\setminus C)\cap \J_h\bigr)\) \\
\bottomrule
\end{tabular}
\end{center}
\end{proposition}

\begin{proof}[Proof sketch]
Each formula is a direct rewriting of the gate semantics in terms of bit conditions over the connected inputs. Monotone gates reduce to intersections and unions of one-bands; parity gates reduce to odd/even weight classes; threshold gates reduce to Hamming-weight strata; implication gates reduce to the single forbidden pattern \(a=1,b=0\); and canalising gates reduce to a collapse on the trigger set plus fallback behaviour on its complement.
\end{proof}

\begin{remark}
The decisive point is that these are \emph{network-aware} formulae. The local rule is not written abstractly over unnamed inputs; it is anchored to the actual connected-input set \(\Ic\) induced by the network topology.
\end{remark}

\subsection{Worked Micro-Examples}
Two small examples already reveal the method's logic.

\paragraph{AND on \(\Ic=\{2,3\}\) with \(n=3\).}
Under LSB-first ordering, node \(2\) contributes the band \(\Band{2}{1}=\{3,4,7,8\}\) and node \(3\) contributes \(\Band{3}{1}=\{5,6,7,8\}\). Their intersection is
\[
\J^{\mathrm{AND}}=\{7,8\}\cap\{j\in\U : v_2(j)=v_3(j)=1\}.
\]
When the explicit ordering convention is enforced against the repertoire generator, the valid index set for the connected bits is the unique all-ones connected pattern. The important methodological lesson is not the accidental numeric value of the index but the need to fix the ordering and hold it fixed throughout derivation and validation.

\paragraph{OR on \(\Ic=\{2,3\}\) with \(n=3\).}
The OR one-set is
\[
\J^{\mathrm{OR}}=\Band{2}{1}\cup\Band{3}{1}=\{3,4,5,6,7,8\},
\]
exactly matching the empirical ordered repertoire when the same ordering policy is used on both sides.

\section{Ordering Invariance}
The formalism would be scientifically brittle if its exactness depended on accidental bit-order conventions. The next theorem shows that it does not.

\begin{theorem}[Ordering-Transform Invariance]
Let \(\J^{\mathrm{LSB}}\subseteq\U\) be the one-set of any gate family under LSB-first ordering, and let \(\J^{\mathrm{MSB}}\) be the corresponding one-set under MSB-first ordering. Then
\[
\J^{\mathrm{MSB}}=\varphi\bigl(\J^{\mathrm{LSB}}\bigr).
\]
\end{theorem}

\begin{proof}
Bit reversal exchanges the LSB-first and MSB-first encodings of the same binary vector. Every gate family in Proposition 1 is defined by conditions on bit values, parity, or Hamming weight over the connected inputs. These conditions are invariant under relabelling of the coordinate positions when the relabelling is carried consistently through the indexing map. Therefore the numeric index sets differ only by transport through \(\varphi\).
\end{proof}

\begin{corollary}
All correctness claims in this paper are representation-invariant provided that the declared ordering convention is used consistently in both derivation and validation.
\end{corollary}

This theorem is central to the paper's causal stance. It isolates the mechanism from the encoding. Two enumerations may look numerically different while representing the same causal rule.

\section{Network Composition and Exact Reconstruction}
The local formulas above become scientifically valuable only if they compose into exact network-level predictions.

\begin{definition}[Analytic Predictor]
For each node \(i\), let \(\J_i\subseteq\U\) denote the one-set determined by its gate family, parameters, connected-input set \(N_i\), and the declared ordering policy. The analytic predictor reconstructs the network output matrix \(Y\in\bits^{2^n\times n}\) by declaring
\[
Y_{j,i}=1 \iff j\in \J_i.
\]
\end{definition}

\begin{theorem}[Canonical Exact Reconstruction]
Fix a network \((cm,dynamic)\), its gate parameters, and an ordering policy. Then the analytic predictor coincides exactly with the exhaustive one-step synchronous baseline:
\[
Y^{\mathrm{analytic}}=Y^{\mathrm{baseline}}.
\]
\end{theorem}

\begin{proof}[Proof sketch]
For each node \(i\), Proposition 1 gives a one-set \(\J_i\) equivalent to the local gate semantics applied to the connected inputs \(x_{N_i}\). Therefore column \(i\) of the analytic output matrix equals column \(i\) of the baseline matrix for every repertoire index \(j\). Equality of all columns implies equality of the matrices.
\end{proof}

\begin{remark}
This theorem does not say that the method avoids the existence of \(2^n\) states. It says something stronger and more relevant: the full network behaviour can be reconstructed by causal formulas rather than by opaque row-wise recomputation.
\end{remark}

\section{Complexity Analysis}
High-level journal quality requires more than exactness. It requires a clear account of what is computationally improved and what is not.

\subsection{What Cannot Be Avoided}
Any method that explicitly outputs the full repertoire matrix of an \(n\)-node network must produce \(n2^n\) bits, hence has an output-size lower bound of \(\Omega(n2^n)\). No honest analysis should claim otherwise.

\subsection{What the Method Improves}
The contribution lies in replacing brute-force local evaluation by exact causal structure.

\begin{proposition}[Per-Row Cost Comparison]
Let the in-degree of node \(i\) be \(|N_i|\). For a fixed input row \(x\),
\begin{itemize}
\item predictive gate-semantic evaluation has cost \(\Theta\!\bigl(\sum_i |N_i|\bigr)\) when local rules are applied directly to the connected inputs;
\item truth-table lookup requires \(\Theta\!\bigl(\sum_i 2^{|N_i|}\bigr)\) preprocessing to construct tables, amortised under reuse;
\item analytic index-set membership uses closed forms built from band, parity, or threshold predicates, again with cost bounded by the connected-input structure rather than by brute-force enumeration of local truth tables.
\end{itemize}
\end{proposition}

\begin{proposition}[Validation-Burden Reduction]
Exhaustive validation requires \(\Theta(2^n)\) rows. Stratified importance sampling validates equality on a fixed \(m\) sampled rows, preserving the exact method while reducing the validation burden from exponential in \(n\) to linear in \(m\).
\end{proposition}

The crucial distinction is therefore between \emph{method complexity} and \emph{output complexity}. The method remains exact. Sampling enters only as a way of auditing exactness at scales where full baseline enumeration becomes computationally undesirable.

\subsection{Concrete Empirical Comparison}
The repository already contains a direct mixed-gate comparison on a 10-node network.
\begin{center}
\begin{tabular}{p{0.29\linewidth}p{0.12\linewidth}p{0.18\linewidth}p{0.23\linewidth}}
\toprule
\textbf{Path} & \textbf{Accuracy} & \textbf{Time} & \textbf{Status} \\
\midrule
Vectorised shortcut formula path & \(0.66875\) & \(1.10\times 10^{-3}\,\mathrm{s}\) & fast but not exact \\
Exact library semantics & \(1.0\) & \(0.081\,\mathrm{s}\) & exact \\
Analytic index-set formulae & \(1.0\) & \(9.85\times 10^{-3}\,\mathrm{s}\) & exact and fastest exact path \\
Baseline exhaustive generation & -- & \(0.118\,\mathrm{s}\) & reference \\
\bottomrule
\end{tabular}
\end{center}

This table matters for two reasons. First, it demonstrates that naive speed alone is not a scientific victory: the fastest non-exact shortcut is unusable as a theorem-supporting method. Second, it shows that analytic formulas achieve what a serious causality-based method should achieve: exactness plus a meaningful computational advantage over other exact baselines.

\section{Validation Evidence}
The current repository supports a substantial validation base. The results most relevant to the main manuscript are summarised below.

\begin{center}
\begin{tabular}{p{0.25\linewidth}p{0.12\linewidth}p{0.12\linewidth}p{0.34\linewidth}}
\toprule
\textbf{Regime} & \textbf{Size} & \textbf{Accuracy} & \textbf{Meaning} \\
\midrule
Medium-scale exact reconstruction & \(n=10\) & \(1.0\) & Predictive reconstruction equals exhaustive baseline \\
Medium-scale exact reconstruction & \(n=13\) & \(1.0\) & Predictive reconstruction equals exhaustive baseline \\
Large-scale sampled validation & \(n=20\), \(1020\) samples & \(1.0\) & Exact equality on stratified sampled rows \\
Large-scale sampled validation & \(n=50\), \(1020\) samples & \(1.0\) & Exact equality on stratified sampled rows \\
Per-gate analysis suite & AND, OR, XOR, XNOR, NAND, NOR, NOT, IMPLIES, KOFN, CANALISING & exact & Gate-level formulae and ordering checks pass \\
\bottomrule
\end{tabular}
\end{center}

The gate-analysis layer is especially important. It verifies not just end-to-end equality but the mathematical pieces that justify the reconstruction theorem: ordering invariance, closure under set operations, parity-class behaviour, threshold strictness, and canalising collapse.

\section{A Structured Example from Mixed Dynamics}
The thesis-era analyses remain valuable because they show how apparently irregular behaviour becomes intelligible when node outputs are isolated and the repertoire ordering is respected. Consider the 7-node mixed network defined by
\[
cm_{07}=
\begin{bmatrix}
0&0&1&0&0&0&1\\
0&0&1&0&0&1&0\\
1&0&0&0&1&0&1\\
1&0&1&0&1&0&1\\
0&0&1&1&0&1&1\\
1&1&1&0&0&0&0\\
0&1&0&1&1&1&0
\end{bmatrix},
\qquad
dynamic_{07}=(\mathrm{AND},\mathrm{OR},\mathrm{OR},\mathrm{AND},\mathrm{OR},\mathrm{OR},\mathrm{AND}).
\]

From a raw repertoire table, the output may look combinatorially messy. But once isolated node-by-node, repeated structure appears. This is precisely the methodological point: complexity in the full table does not preclude simple local rules. On the contrary, the recovery of local causal formulas from apparently rich behaviour is one of the main reasons the method is scientifically interesting.

This perspective places the method near the broader insight that simple local rules can generate surprisingly rich global behaviour. The present work differs, however, in that it does not stop at the philosophical observation. It provides explicit formulas for locating that behaviour within the repertoire.

\section{Relation to Algorithmic Causality and Statistical Learning}
\subsection{Relation to Wolfram-Style Rule Simplicity}
Wolfram's work established that simple discrete rules can generate unexpectedly rich behaviour. Our contribution is compatible with that viewpoint but more operational for Boolean networks: we derive exact formulas that specify where rule-induced behaviour appears in an ordered repertoire, rather than merely demonstrating that rich behaviour exists.

\subsection{Relation to Algorithmic Causality}
Algorithmic-causality programmes emphasise perturbation, generative explanation, and the distinction between statistical regularity and genuine causal mechanism. The present method fits naturally into that tradition because it replaces descriptive output summaries by explicit program-like rule expressions. In this sense, the index-set calculus is a concrete causal representation.

\subsection{Difference from Classical Machine Learning}
The contrast with classical machine learning should be stated carefully. Many ML systems learn predictive associations from sampled data and are indifferent to whether the learned representation coincides with the true generating rule. Here the objective is different: under known network logic, the goal is to derive exact consequences analytically. The paper is therefore not anti-statistical; it is simply aimed at a different problem class, one in which causal structure is explicit and deterministic reconstruction is possible.

\section{Discussion}
For a high-level journal, the central value of the method is not that it is ``faster than brute force'' in some vague sense. The stronger claim is that it changes the scientific unit of explanation. The unit is no longer the truth table or the trajectory, but the exact rule-induced index set over an ordered repertoire.

This has four consequences.
\begin{enumerate}[label=\arabic*.]
\item It yields \emph{explanation}: the one-set specifies why a node is active and not merely when it happened to be active in a listed table.
\item It yields \emph{representation control}: ordering changes are transported by \(\varphi\) rather than handled by undocumented convention changes.
\item It yields \emph{compositionality}: mixed networks inherit their structure from local gate formulas.
\item It yields \emph{computational discipline}: approximate shortcuts are separated cleanly from exact analytic prediction.
\end{enumerate}

The method should therefore be read as a causal-computational mathematics paper with later applications to biological and other complex systems, not as a biology paper wearing mathematical vocabulary. That distinction matters for publication strategy. The main paper must foreground exact definitions, theorems, worked examples, and exact validation. Broader applications can then be layered onto a firmer foundation.

\section{Conclusion}
We have presented a causality-first index-set calculus for Boolean networks. The formalism begins with ordered exhaustive repertoires, defines exact gate-family one-sets, controls ordering by an explicit bit-reversal involution, and composes local one-sets into exact network-level reconstruction. The resulting framework is mathematically cleaner than brute-force enumeration, more explanatory than raw truth-table lookup, and more appropriate than probabilistic surrogates when the underlying mechanism is known.

The next stage is not to reinvent the method but to deepen this manuscript into full journal form: expand the appendix-level derivations for every gate family, insert the strongest result tables and figures from the validated artefacts, and connect the present exact formalism to the downstream paper programme. The core scientific claim, however, is already clear: complex Boolean-network behaviour can be represented and reconstructed by short exact causal formulas once the repertoire ordering is treated as part of the mathematics rather than as an implementation detail.

\appendix

\section{Appendix A: Exact Gate Families in Expanded Form}
\subsection{Monotone Gates}
For a connected-input set \(\Ic\),
\[
\J^{\mathrm{AND}}=\bigcap_{i\in\Ic}\Band{i}{1},
\qquad
\J^{\mathrm{OR}}=\bigcup_{i\in\Ic}\Band{i}{1},
\]
\[
\J^{\mathrm{NAND}}=\U\setminus\J^{\mathrm{AND}},
\qquad
\J^{\mathrm{NOR}}=\U\setminus\J^{\mathrm{OR}}=\bigcap_{i\in\Ic}\Band{i}{0}.
\]

\subsection{Parity and Equivalence}
\[
\J^{\mathrm{XOR}}=\Parity{\Ic}{1},
\qquad
\J^{\mathrm{XNOR}}=\Parity{\Ic}{0}.
\]
These classes partition \(\U\) and are exact complements of one another.

\subsection{Implication Families}
For a designated ordered pair \((a,b)\),
\[
\J^{\mathrm{IMP}}=\U\setminus\bigl(\Band{a}{1}\cap\Band{b}{0}\bigr),
\qquad
\J^{\mathrm{NIMP}}=\Band{a}{1}\cap\Band{b}{0}.
\]
The asymmetry is essential and must not be erased by symmetrising the pair.

\subsection{Threshold Families}
For KOFN with threshold \(k\),
\[
\J^{\mathrm{KOFN}}=\bigcup_{r=k}^{|\Ic|}\Weight{\Ic}{r}.
\]
For exact-\(k\) thresholds,
\[
\J^{=k}=\Weight{\Ic}{k}.
\]
The special cases \(k=0\) and \(k>|\Ic|\) yield the constant functions \(1\) and \(0\), respectively.

\subsection{Canalising Families}
Let \(C\subseteq\U\) be the trigger set on which the canalising value is detected and let \(\J_h\) be the one-set of the fallback rule. Then
\[
\J^{\mathrm{CAN}}=C\cup\bigl((\U\setminus C)\cap \J_h\bigr).
\]
Nested canalising rules iterate this construction with priority-ordered trigger sets.

\section{Appendix B: Sensitivity and Interpretive Notes}
The exact gate formulas support additional analyses already present in the repository, including average sensitivity and structured noise response.

\paragraph{AND and OR.}
For arity \(d\),
\[
\bar{s}_{\mathrm{AND}}=\bar{s}_{\mathrm{OR}}=\frac{d}{2^{d-1}}.
\]

\paragraph{XOR and XNOR.}
For arity \(d\),
\[
\bar{s}_{\mathrm{XOR}}=\bar{s}_{\mathrm{XNOR}}=d,
\]
since flipping any single input toggles parity.

\paragraph{Interpretive point.}
Sensitivity is not the primary theorem of this paper, but it matters because it connects the static one-set calculus to perturbation behaviour. The same exact formulas that identify where outputs equal one also determine how readily those outputs change under controlled input perturbations.

\section{Appendix C: Validation Summary and Reproducibility}
The manuscript is grounded in a larger validation programme covering:
\begin{itemize}
\item ordering invariance checks under \(\varphi\),
\item exact gate-level one-set equality against empirical repertoires,
\item medium-scale exact reconstruction for \(n=10,13\),
\item large-scale sampled equality for \(n=20,50\),
\item mixed-gate comparisons separating exact analytic methods from intentionally inexact shortcut methods.
\end{itemize}

The main artefact families include \texttt{results/tests/analysis\_*}, \texttt{results/tests/algo001}, \texttt{results/tests/algo002}, and \texttt{results/tests/mixed001FormulaVsExhaustive}. These artefacts are important not as supplementary clutter but as audit trails for every main claim made here.

\begin{thebibliography}{99}
\bibitem{wolfram}
S. Wolfram, \emph{A New Kind of Science}, Wolfram Media, 2002.

\bibitem{zenilcausality}
H. Zenil, N. A. Kiani, and J. Tegner, ``Low-algorithmic-complexity entropy-deceiving graphs,'' \emph{Physical Review E}, 96(1), 2017.

\bibitem{zenildynamics}
H. Zenil, F. Soler-Toscano, N. A. Kiani, S. Hern\'andez-Orozco, and J. Tegn\'er, ``A decomposition method for global evaluation of Shannon entropy and local estimations of algorithmic complexity,'' \emph{Entropy}, 20(8), 2018.

\bibitem{li}
M. Li and P. Vit\'anyi, \emph{An Introduction to Kolmogorov Complexity and Its Applications}, 4th ed., Springer, 2019.

\bibitem{pearl}
J. Pearl, \emph{Causality: Models, Reasoning, and Inference}, 2nd ed., Cambridge University Press, 2009.

\bibitem{kauffman}
S. A. Kauffman, ``Metabolic stability and epigenesis in randomly constructed genetic nets,'' \emph{Journal of Theoretical Biology}, 22(3), 1969.
\end{thebibliography}

\end{document}
